\section{Results and Discussion}
\label{sec:results}

This section presents the empirical findings of our comparative study, organized into four subsections addressing convergence behavior, dimensional scaling, quantum mode differences, and statistical significance. We conclude with a broader discussion of the practical implications.

\subsection{Convergence Analysis}
\label{subsec:convergence}

Figure~\ref{fig:convergence-curves} presents the convergence curves for Classical PSO and QPSO across all six benchmark functions at representative dimensionalities. Each curve depicts the median global best fitness over 30 independent runs, with the shaded region indicating the interquartile range.

On the Sphere function, both algorithms converge rapidly toward the global optimum, with QPSO achieving a final best fitness of \textbf{[TBD]} compared to Classical PSO's \textbf{[TBD]} in the 10-dimensional case. The convergence trajectories are largely overlapping during the first \textbf{[TBD]} iterations, after which QPSO exhibits a marginally steeper descent. This behavior is expected, as the smooth, convex landscape of the Sphere function does not present the local optima that would differentiate the exploration mechanisms of the two algorithms.

The multimodal benchmark functions reveal more pronounced differences. On Rastrigin, QPSO achieves a median final fitness of \textbf{[TBD]}, representing a \textbf{[TBD]}\% improvement over Classical PSO's median of \textbf{[TBD]} in 10 dimensions. The convergence curves show that Classical PSO frequently stagnates at local optima after approximately \textbf{[TBD]} iterations, whereas QPSO's quantum-inspired position update mechanism enables continued exploration of the search space. A similar pattern is observed on the Ackley function, where QPSO's final median fitness of \textbf{[TBD]} compares favorably with Classical PSO's \textbf{[TBD]}.

The Rosenbrock function provides perhaps the most interesting comparison. The narrow, curved valley topology rewards algorithms capable of making incremental adjustments along the ridge while maintaining sufficient momentum to traverse its curvature. QPSO with the weighted mean best position achieves a median final fitness of \textbf{[TBD]}, compared to \textbf{[TBD]} for Classical PSO. The convergence curves in Figure~\ref{fig:convergence-rosenbrock} show that QPSO's advantage emerges primarily in the later stages of the optimization, suggesting that the contraction--expansion coefficient's decay schedule facilitates a smooth transition from exploration to fine-grained exploitation along the valley floor.

On Griewank, both algorithms reach near-optimal solutions in low dimensions, with final fitness values of \textbf{[TBD]} (QPSO) and \textbf{[TBD]} (PSO) in the 5-dimensional case. However, as dimensionality increases, the performance gap widens, reflecting the compounding effect of the product term in the Griewank function across higher-dimensional spaces.

The Schwefel function, with its deceptive landscape and geometrically distant second-best minimum, presents the greatest challenge. Neither algorithm consistently reaches the global optimum in higher dimensions, but QPSO achieves a lower median residual error of \textbf{[TBD]} compared to \textbf{[TBD]} for Classical PSO at $n = 20$. The box plots in Figure~\ref{fig:boxplots-schwefel} reveal that Classical PSO's distribution of final fitness values is notably wider, indicating less consistent performance across independent runs.

\subsection{Dimensional Scaling}
\label{subsec:scaling}

To assess how each algorithm copes with the curse of dimensionality, we examine performance trends across the four tested dimensionalities: $n \in \{2, 5, 10, 20\}$. Table~\ref{tab:scaling-results} reports the mean and standard deviation of the final best fitness for each algorithm--benchmark--dimensionality combination.

\begin{table}[htbp]
\centering
\caption{Mean $\pm$ standard deviation of final best fitness across 30 independent runs. Bold entries indicate the better-performing algorithm for each benchmark--dimensionality pair.}
\label{tab:scaling-results}
\small
\begin{tabular}{@{}llcccc@{}}
\toprule
\textbf{Function} & \textbf{Algorithm} & $n=2$ & $n=5$ & $n=10$ & $n=20$ \\
\midrule
\multirow{2}{*}{Sphere}
 & PSO  & \textbf{[TBD]} & \textbf{[TBD]} & \textbf{[TBD]} & \textbf{[TBD]} \\
 & QPSO & \textbf{[TBD]} & \textbf{[TBD]} & \textbf{[TBD]} & \textbf{[TBD]} \\
\midrule
\multirow{2}{*}{Rastrigin}
 & PSO  & \textbf{[TBD]} & \textbf{[TBD]} & \textbf{[TBD]} & \textbf{[TBD]} \\
 & QPSO & \textbf{[TBD]} & \textbf{[TBD]} & \textbf{[TBD]} & \textbf{[TBD]} \\
\midrule
\multirow{2}{*}{Rosenbrock}
 & PSO  & \textbf{[TBD]} & \textbf{[TBD]} & \textbf{[TBD]} & \textbf{[TBD]} \\
 & QPSO & \textbf{[TBD]} & \textbf{[TBD]} & \textbf{[TBD]} & \textbf{[TBD]} \\
\midrule
\multirow{2}{*}{Ackley}
 & PSO  & \textbf{[TBD]} & \textbf{[TBD]} & \textbf{[TBD]} & \textbf{[TBD]} \\
 & QPSO & \textbf{[TBD]} & \textbf{[TBD]} & \textbf{[TBD]} & \textbf{[TBD]} \\
\midrule
\multirow{2}{*}{Griewank}
 & PSO  & \textbf{[TBD]} & \textbf{[TBD]} & \textbf{[TBD]} & \textbf{[TBD]} \\
 & QPSO & \textbf{[TBD]} & \textbf{[TBD]} & \textbf{[TBD]} & \textbf{[TBD]} \\
\midrule
\multirow{2}{*}{Schwefel}
 & PSO  & \textbf{[TBD]} & \textbf{[TBD]} & \textbf{[TBD]} & \textbf{[TBD]} \\
 & QPSO & \textbf{[TBD]} & \textbf{[TBD]} & \textbf{[TBD]} & \textbf{[TBD]} \\
\bottomrule
\end{tabular}
\end{table}

In the low-dimensional setting ($n = 2$), both algorithms perform comparably across all benchmarks, with neither showing a consistent advantage. This is consistent with the observation that two-dimensional landscapes are generally tractable for swarm-based methods, as the search space volume remains small enough for adequate coverage with 40 particles.

As dimensionality increases, the performance gap between the two algorithms becomes more pronounced on multimodal functions. On Rastrigin at $n = 20$, the number of local minima scales as $10^{20}$, and Classical PSO's velocity-based exploration mechanism becomes increasingly insufficient for navigating such a rugged landscape. QPSO's position update equation, which samples from a probability distribution centered around the weighted mean best, provides a fundamentally different exploration mechanism that appears better suited to high-dimensional multimodal search.

On the unimodal Sphere function, the dimensional scaling effect is less dramatic. Both algorithms maintain good convergence properties even at $n = 20$, though the absolute fitness values increase with dimensionality, reflecting the inherent difficulty of high-dimensional optimization. The relative performance difference remains modest, confirming that the quantum-inspired update offers limited incremental benefit on smooth landscapes.

The Rosenbrock function exhibits an interesting scaling pattern. At $n = 2$, the function has a single minimum, and both algorithms converge reliably. At higher dimensions, the function becomes increasingly challenging due to the sequential dependence among variables, with the valley narrowing and lengthening as $n$ grows. QPSO's advantage on this function becomes more pronounced at $n = 10$ and $n = 20$, suggesting that the weighted mean best mechanism facilitates coordinated movement along the valley ridge.

\subsection{Quantum Mode Comparison}
\label{subsec:quantum-modes}

A distinctive feature of our framework is the support for three quantum randomness generation modes within QPSO: mathematical simulation (Math), full Qiskit quantum circuit execution (Full), and a hybrid approach (Hybrid) that uses quantum circuits for a subset of the random number generation while relying on classical pseudorandom numbers for the remainder.

Table~\ref{tab:quantum-modes} presents the median final fitness and mean wall-clock time per run for each mode across all benchmarks at $n = 10$.

\begin{table}[htbp]
\centering
\caption{Comparison of QPSO quantum randomness modes at $n = 10$. Median final fitness (lower is better) and mean wall-clock time per run are reported across 30 independent runs.}
\label{tab:quantum-modes}
\small
\begin{tabular}{@{}lccc@{}}
\toprule
\textbf{Function} & \textbf{Math Mode} & \textbf{Full Mode} & \textbf{Hybrid Mode} \\
\midrule
Sphere    & \textbf{[TBD]} / \textbf{[TBD]}s & \textbf{[TBD]} / \textbf{[TBD]}s & \textbf{[TBD]} / \textbf{[TBD]}s \\
Rastrigin & \textbf{[TBD]} / \textbf{[TBD]}s & \textbf{[TBD]} / \textbf{[TBD]}s & \textbf{[TBD]} / \textbf{[TBD]}s \\
Rosenbrock& \textbf{[TBD]} / \textbf{[TBD]}s & \textbf{[TBD]} / \textbf{[TBD]}s & \textbf{[TBD]} / \textbf{[TBD]}s \\
Ackley    & \textbf{[TBD]} / \textbf{[TBD]}s & \textbf{[TBD]} / \textbf{[TBD]}s & \textbf{[TBD]} / \textbf{[TBD]}s \\
Griewank  & \textbf{[TBD]} / \textbf{[TBD]}s & \textbf{[TBD]} / \textbf{[TBD]}s & \textbf{[TBD]} / \textbf{[TBD]}s \\
Schwefel  & \textbf{[TBD]} / \textbf{[TBD]}s & \textbf{[TBD]} / \textbf{[TBD]}s & \textbf{[TBD]} / \textbf{[TBD]}s \\
\bottomrule
\end{tabular}
\end{table}

In terms of solution quality, the three modes produce largely comparable final fitness values. The Math mode, which generates quantum-distributed random numbers using classical transformations, achieves median fitness values within \textbf{[TBD]}\% of the Full mode across all benchmarks. This suggests that the mathematical approximation of quantum randomness is sufficient for the optimization dynamics, and the precise source of stochasticity has a secondary effect compared to the structural differences in the position update equation.

The computational overhead, however, differs substantially across modes. The Full mode, which constructs and simulates Qiskit quantum circuits with 256 shots per random number generation, incurs a wall-clock time overhead of approximately \textbf{[TBD]}$\times$ compared to the Math mode. This overhead arises from circuit construction, transpilation, and statevector simulation on a classical backend. The Hybrid mode occupies a middle ground, with a time overhead of approximately \textbf{[TBD]}$\times$, as it delegates only a fraction of the random number generation to the quantum circuit pathway.

These findings have practical implications for the deployment of QPSO. For researchers interested in studying the algorithmic behavior of QPSO without access to quantum hardware, the Math mode provides an efficient and faithful approximation. The Full and Hybrid modes become more relevant in contexts where the true quantum randomness properties---such as provable unpredictability and uniform distribution guarantees from quantum measurement---are of intrinsic interest, or when the framework is adapted for execution on real quantum processing units.

\subsection{Statistical Significance}
\label{subsec:significance}

Table~\ref{tab:wilcoxon} reports the results of the Wilcoxon signed-rank test comparing Classical PSO and QPSO (Math mode) across all benchmark--dimensionality combinations. The test is applied to the paired final best fitness values from the 30 independent runs.

\begin{table}[htbp]
\centering
\caption{Wilcoxon signed-rank test results comparing Classical PSO vs.\ QPSO. The winner column indicates which algorithm achieved statistically significantly lower median fitness ($\alpha = 0.05$). A dash indicates no statistically significant difference.}
\label{tab:wilcoxon}
\small
\begin{tabular}{@{}lcccc@{}}
\toprule
\textbf{Function} & $n=2$ & $n=5$ & $n=10$ & $n=20$ \\
\midrule
Sphere     & \textbf{[TBD]} & \textbf{[TBD]} & \textbf{[TBD]} & \textbf{[TBD]} \\
Rastrigin  & \textbf{[TBD]} & \textbf{[TBD]} & \textbf{[TBD]} & \textbf{[TBD]} \\
Rosenbrock & \textbf{[TBD]} & \textbf{[TBD]} & \textbf{[TBD]} & \textbf{[TBD]} \\
Ackley     & \textbf{[TBD]} & \textbf{[TBD]} & \textbf{[TBD]} & \textbf{[TBD]} \\
Griewank   & \textbf{[TBD]} & \textbf{[TBD]} & \textbf{[TBD]} & \textbf{[TBD]} \\
Schwefel   & \textbf{[TBD]} & \textbf{[TBD]} & \textbf{[TBD]} & \textbf{[TBD]} \\
\bottomrule
\end{tabular}
\end{table}

Across the \textbf{[TBD]} benchmark--dimensionality combinations tested, QPSO achieves a statistically significant advantage in \textbf{[TBD]} cases, Classical PSO in \textbf{[TBD]} cases, and no significant difference is observed in the remaining \textbf{[TBD]} cases. The pattern confirms the qualitative observations from the convergence analysis: QPSO's advantages concentrate on multimodal and valley-shaped functions at higher dimensionalities, while the algorithms are statistically indistinguishable on the Sphere function in low dimensions.

The Friedman test, applied across all four algorithm variants (Classical PSO, QPSO-Math, QPSO-Full, QPSO-Hybrid), yields a test statistic of $\chi^2_F = $ \textbf{[TBD]} with a $p$-value of \textbf{[TBD]}, indicating a statistically significant difference in overall rankings. The mean ranks are \textbf{[TBD]} for Classical PSO, \textbf{[TBD]} for QPSO-Math, \textbf{[TBD]} for QPSO-Full, and \textbf{[TBD]} for QPSO-Hybrid. Post-hoc analysis with Holm correction reveals that \textbf{[TBD]}.

\subsection{Discussion}
\label{subsec:discussion}

The empirical evidence presented in this study supports several conclusions regarding the relative merits of Classical PSO and Quantum-behaved PSO.

QPSO demonstrates its clearest advantages on functions characterized by multimodality and narrow-valley topologies. On the Rastrigin and Ackley functions, the quantum-inspired position update equation provides a qualitatively different exploration mechanism compared to the velocity-based movement of Classical PSO. While Classical PSO's particle velocities tend to decay as the swarm converges, potentially trapping particles in the basin of attraction of a local optimum, QPSO's probabilistic sampling from a distribution centered around the mean best position allows particles to ``tunnel'' through potential barriers in a manner conceptually analogous to quantum tunneling. This effect is most pronounced in higher dimensions, where the density of local optima increases exponentially.

The weighted mean best position, computed as a fitness-proportional average of all personal best positions in the swarm, plays a particularly important role on the Rosenbrock function. Unlike the simple arithmetic mean, the weighted variant biases the attractor point toward the personal bests of higher-quality particles, effectively guiding the swarm along the narrow valley floor. This mechanism complements the contraction--expansion coefficient's decay schedule, which narrows the sampling distribution as the optimization progresses, enabling progressively finer adjustments near the optimum.

The stagnation detection and recovery mechanism, shared by both algorithms, contributes to the robustness of the results. By perturbing a fraction of the swarm when no improvement is observed for 15 consecutive iterations, the mechanism prevents complete convergence to a suboptimal solution. However, the effectiveness of this mechanism differs between the two algorithms: in Classical PSO, stagnation recovery relies on velocity perturbation, which may not generate sufficient displacement in a high-dimensional space, while in QPSO, the inherent stochasticity of the position update provides a complementary source of diversity that enhances the recovery mechanism's effectiveness.

From a practical standpoint, the computational cost of the quantum circuit modes warrants careful consideration. The Full Qiskit mode, while faithful to the quantum-mechanical formulation, incurs significant overhead due to circuit simulation on classical hardware. For most optimization applications where solution quality is the primary concern, the Math mode provides an efficient alternative with negligible loss in performance. The Full and Hybrid modes are better positioned for future deployment on actual quantum hardware, where the circuit execution overhead would be replaced by the physical operation of quantum gates, potentially offering speedups for high-dimensional problems that require large numbers of random samples.
